%% v74_amendment.tex
%% "ECI v7.4 amendment: tau as a CM-anchored attractor and a revised
%%  Damerell-ladder bridge H7' for Cardy density at c = 1 and c = 1/2"
%% Target: Letters in Mathematical Physics (LMP) — short note, 8-12 pp
%% Draft date: 2026-05-05 (afternoon, post G1.15 audit)
%%
%% Anti-hallucination policy:
%%   - All arXiv IDs cited below are live-verified (arXiv API, 2026-05-05).
%%   - Hallucination counter held at 76; do NOT introduce new fabricated cites.
%%   - The Damerell ladder (1/10, 1/12, 1/24, 1/60) is the standard
%%     Eisenstein-Kronecker ladder for a weight-5 CM-by-Q(i) newform
%%     (Hurwitz 1899 + Chowla-Selberg 1949 + Damerell 1971 + Shimura 1976
%%     + Harder-Schappacher 1986). Verified literature anchor (A1, 2026-05-05).

\documentclass[a4paper,11pt]{article}

\usepackage{amsmath,amssymb,amsfonts,amsthm}
\usepackage{hyperref}
\usepackage{xcolor}
\usepackage{booktabs}
\usepackage{geometry}
\usepackage{bm}
\usepackage{url}
\geometry{a4paper, margin=2.5cm}

%% Shorthand macros (consistent with V2 paper)
\newcommand{\Spf}{S'_4}
\newcommand{\tauCM}{\tau = i}
\newcommand{\arXiv}[1]{\href{https://arxiv.org/abs/#1}{\texttt{arXiv:#1}}}
\newcommand{\hhat}[1]{\widehat{\mathbf{#1}}}
\newcommand{\OmK}{\Omega_K}
\newcommand{\Lf}{L(f,\,\cdot\,)}
\newcommand{\fLM}{f_{\text{4.5.b.a}}}
\newcommand{\Qi}{\mathbb{Q}(i)}
\newcommand{\Zeno}{\href{https://doi.org/10.5281/zenodo.20030684}{\texttt{10.5281/zenodo.20030684}}}

\newtheorem{theorem}{Theorem}
\newtheorem{proposition}[theorem]{Proposition}
\newtheorem{lemma}[theorem]{Lemma}
\newtheorem{conjecture}[theorem]{Conjecture}
\newtheorem*{axiom*}{Axiom}
\theoremstyle{definition}
\newtheorem{definition}[theorem]{Definition}
\theoremstyle{remark}
\newtheorem{remark}[theorem]{Remark}

\title{ECI v7.4 amendment:\\
       $\tau$ as a CM-anchored attractor and a revised\\
       Damerell-ladder bridge $H_7'$ for Cardy density\\
       at $c=1$ and $c=1/2$}

\author{%
  K.~Remondi\`ere%
  \\[2ex]
  \small\textit{Independent researcher, Tarbes, France}\\
  \small\texttt{kevin.remondiere@gmail.com}
}

\date{Draft: 2026-05-05}

\begin{document}
\maketitle

%%--------------------------------------------------------------------
\begin{abstract}
The v7.3 axiomatic system of the Extended Cosmological Index (ECI)
research programme was archived in version~6.0.53 of the project
repository (Zenodo DOI \Zeno) on 2026-05-05.  An internal Lakatos-style
audit (G1.15) identified one refuted axiom: $H_7$ asserted a
Damerell--Coates--Sinnott bridge for the LMFDB CM newform
$\fLM$ (CM by $\Qi$, weight $k=5$) at the central
half-integer $s=k/2=5/2$, which lies \emph{outside} the
Damerell/Shimura/Deligne algebraicity domain for odd~$k$.
\par
This note records the v7.3 $\to$ v7.4 amendment.  Three independent
sub-agents (H7-A, H7-B, H7-C) explored rescue paths and converged on
keeping $\fLM$ as the v7.4 CM anchor.  We propose two replacements.
First, $H_5'$ relaxes ``$\tau$ is the strict $S$-fixed point $i$'' to
``$\tau$ is a CM-anchored \emph{attractor}, $|\tau-i|\lesssim 0.2$'',
supported empirically by a $30\times 30$ $\chi^2$ scan over seven
fermion observables (W1, PC compute, 2026-05-05) which locates a basin
at $\tau^{*} = -0.1897 + 1.0034\,i$ with $\chi^2/\mathrm{dof}=1.05$ and
$|\tau^{*}-i|=0.19$ -- a factor $\approx\!3$ closer to $i$ than the
LYD20 published best fit~\cite{LYD20}.  Second, $H_7'$ replaces the
half-integer Damerell--CS bridge by (i) an INTEGER-point Damerell
ladder
\[
\boxed{\;\;L(\fLM, m)\,\pi^{4-m}/\OmK^{4} \;\in\; \{1/10,\,1/12,\,1/24,\,1/60\}
       \;\;\text{for}\;\; m\in\{1,2,3,4\}\;\;}
\]
which is the standard Eisenstein--Kronecker ladder for a weight-$5$
CM-by-$\Qi$ newform: $\alpha_1 = 1/10$ is the lemniscatic Hurwitz
number $H_4$~\cite{Hurwitz1899,HS86}, $\alpha_2 = 1/12 = B_2/2 =
-\zeta(-1)$ is the second Bernoulli half-value transported to $\Qi$
via Chowla--Selberg~\cite{ChowSel49,Damerell71,Shimura76}, and
$(\alpha_3,\alpha_4) = (\alpha_2/2, \alpha_1/6)$ are forced by the
$\Gamma$-functional equation $L(\fLM,s) \leftrightarrow L(\fLM,5-s)$.
Therefore the apparent ``two parameter-free Cardy hits'' at $c=1$ and
$c=1/2$ reduce to a \emph{single} Bernoulli-anchored coincidence
($B_2/2 = 1/12 = \rho_{\mathrm{Cardy}}(c=1)$) plus its automatic
$\Gamma$-shadow at $c=1/2$.  As complement, (ii) a
Bertolini--Darmon--Prasanna anti-cyclotomic $p$-adic $L$-value at the
original central point $s=k/2$, used when INTEGER-point algebraicity
is unavailable.
We discuss falsifiers (Tonks--Girardeau density, modular flavour
lepton sector, NMC cosmological NULL at Cassini) and flag one open
question: the failure to cleanly fit $c=7/10$ (tricritical Ising) and
$c=4/5$ (tetracritical Potts) under either bridge route.
\end{abstract}

%%--------------------------------------------------------------------
\section{Setup: ECI v7.3 and the G1.15 refutation of $H_7$}
\label{sec:setup}

The Extended Cosmological Index (ECI) research programme proposes that
a single CM newform anchors both flavour-sector mass ratios (via
$\Spf$ modular flavour models~\cite{NPP20,LYD20,dMVP26}) and
two-dimensional CFT Cardy density at minimal-model central charges,
through the algebraic part of its critical $L$-values normalised by a
power of the Chowla--Selberg period~\cite{ChowSel49}.  In the v7.3
axiomatic restart (committed in repository version 6.0.53, Zenodo DOI
\Zeno), the relevant axioms were:
\begin{itemize}
  \item $H_5$: \emph{strict $S$-fixed point}.  $\tau = i$ is the
    physical value of the modular parameter, motivated by maximal
    residual symmetry $\mathbb{Z}_4\subset\mathrm{SL}(2,\mathbb{Z})$.
  \item $H_7$: \emph{Damerell--Coates--Sinnott bridge}.  The Cardy
    density $\rho_{\mathrm{Cardy}}(c)=c/12$ at minimal-model central
    charges is reproduced by the algebraic part of $L(\fLM, k/2)$ with
    $k=5$, $\fLM$ the LMFDB newform 4.5.b.a (CM by $\Qi$).
\end{itemize}

The companion note~\cite{V2NoGo} (V2) had already rendered $H_5$
suspect: it proves a no-go theorem for the Cabibbo angle in LYD20
Model VI at $\tau=i$ (numerical bound $\sin\theta_C \le 0.005$ versus
PDG $0.2253$).  The internal audit G1.15 (May~2026) then tested $H_7$
against the algebraicity hypotheses of Damerell~\cite{Damerell71},
Shimura~\cite{Shimura76}, and Deligne~\cite{Deligne79} for critical
$L$-values of CM modular forms.  These hypotheses require the critical
point to be a non-negative integer (for the standard $L$-function of a
weight-$k$ CM form, the right-of-center critical points are
$s\in\{k-1, k-2,\ldots, \lceil k/2\rceil\}$; the parity of~$k$
controls which half-integer is included).  For $k=5$ odd, $s=k/2=5/2$
is \emph{not} in the algebraicity domain proved by Damerell--Shimura
nor in the period-extension of Deligne.

\begin{proposition}[Refutation of $H_7$, G1.15 audit]
\label{prop:refute}
The half-integer $s=5/2$ is outside the algebraicity domain of
Damerell~\cite{Damerell71} and Shimura~\cite{Shimura76} for a CM
newform of weight $k=5$ over $\Qi$.  Consequently, the assertion
$L(\fLM, 5/2)\,\pi^{3/2}/\OmK^{4}\in\mathbb{Q}$ is not entailed by
any classical algebraicity theorem on which $H_7$ relied.
The numerical value $L(\fLM, 5/2) = 0.5200744676\ldots$ (LMFDB,
mpmath dps$=60$) is consistent with $H_7$ but not implied by any
known theorem; absent such a theorem, $H_7$ is unfalsifiable in its
v7.3 form and must be replaced.
\end{proposition}

\noindent
The audit log is recorded in repository commit 6.0.53; cf.\ Zenodo
DOI \Zeno.

%%--------------------------------------------------------------------
\section{Sub-agent verdicts on the rescue}
\label{sec:rescue}

Three independent sub-agents (Sonnet harness, isolated worktrees) were
tasked to explore rescue paths, all retaining $\fLM$ as the
candidate CM anchor:
\begin{itemize}
  \item \textbf{H7-A (POSITIVE).}  Compute $L(\fLM, m)$ at the
    INTEGER critical points $m\in\{1,2,3,4\}$ via the
    Iwaniec--Kowalski approximate functional equation
    \cite[Thm.~5.3]{IwaKow04} (mpmath dps$=60$, sanity check
    $L(\fLM,5/2)=0.5200744676$ matching LMFDB).  The
    algebraic parts $L(\fLM,m)\,\pi^{4-m}/\OmK^{4}$ form a
    rational ladder; in particular $1/12$ at $m=2$ and $1/24$ at
    $m=3$ \emph{coincide exactly} with the Cardy
    densities $c/12$ at $c=1$ (free boson) and $c=1/2$ (Ising).
  \item \textbf{H7-B (POSITIVE, novel).}  Perlmutter's degree-4
    $L$-function attached to a 2D CFT torus partition
    function~\cite{Perlmutter25} matches the degree of the
    Rankin--Selberg square $L(\fLM \otimes \fLM, s)$, which has
    INTEGER critical points $s\in\{5,\ldots,9\}$ and Petersson norm
    algebraic up to $\pi^{2k-1}=\pi^{9}$.  This suggests an
    alternative bridge $\rho_{\mathrm{Cardy}}\propto
    L(\fLM\otimes\fLM,m_{0})/\langle\fLM,\fLM\rangle$;
    we keep it as a backup since the explicit map is to be invented.
  \item \textbf{H7-C (STRUCTURAL + alternative).}  A structural
    theorem of LYD20 (\cite[Tab.~Level4 MM]{LYD20}, $\S 3$) shows that
    \emph{every} hatted irrep of $\Spf$ carries odd-weight forms,
    enforced by the action of the central element $R\sim -I$ in the
    metaplectic cover $\widetilde{\Gamma}(4)$ as $(-1)^{k}$.  No
    even-weight CM-by-$\Qi$ form can substitute for $\fLM$ inside the
    $\hhat{3}_{,2}^{(5)}$ embedding -- a structural, not enumerative,
    obstruction.  As an alternative central-point bridge, the
    Bertolini--Darmon--Prasanna anti-cyclotomic $p$-adic
    $L$-functions~\cite{BDP13} \emph{do} handle CM-by-imaginary-quadratic
    forms at half-integer central points; they preserve the original
    $s=k/2$ and replace the archimedean assertion by a $p$-adic one.
\end{itemize}

A scan of seven fermion observables (W1, PC compute, 2026-05-05) on a
$30\times 30$ grid in the $\tau$-plane found a local minimum at
\begin{equation}
  \tau^{*} = -0.1897 + 1.0034\,i,\qquad
  |\tau^{*} - i| = 0.19,\qquad
  \chi^{2}/\mathrm{dof} = 1.05,
\end{equation}
which is a factor $\approx 3$ closer to $i$ than the LYD20 published
best fit $\tau_{\mathrm{LYD20}}=-0.21+1.52\,i$ (with
$|\tau_{\mathrm{LYD20}}-i|=0.56$).  This empirically supports
\emph{relaxing} $H_5$ from a strict fixed-point assertion to an
\emph{attractor} statement.  Full synthesis at
\texttt{notes/eci\_v7\_aspiration/H7\_RESCUE/SYNTHESIS\_2026-05-05.md}
and verdict log at
\texttt{notes/eci\_v7\_aspiration/W1\_TAU\_NEAR\_I/results/W1\_verdict\_2026-05-05.log}
in repository commit 6.0.53.

%%--------------------------------------------------------------------
\section{The amended axioms $H_5'$ and $H_7'$}
\label{sec:amended}

\begin{axiom*}[$H_5'$, CM-anchored attractor]
The modular parameter $\tau$ of the $\Spf$ flavour sector is not
fixed at the $S$-fixed point $\tau=i$ but is a \emph{CM-anchored
attractor}: any joint fit to PDG fermion mass ratios and CKM/PMNS
mixings should land within
\[
  |\tau - i| \;\lesssim\; 0.2.
\]
The empirical fit $\tau^{*}=-0.1897+1.0034\,i$ (W1 scan,
$\chi^{2}/\mathrm{dof}=1.05$, 7 observables) is the v7.4 reference
point.
\end{axiom*}

\begin{remark}
The relaxation $H_5\to H_5'$ is forced by V2~\cite{V2NoGo}: at strict
$\tau=i$, LYD20 Model~VI's $\{d^{c},s^{c}\}$-block of the down-quark
matrix becomes near-rank-1 by an $S$-collinearity lemma, and
$\sin\theta_C\le 0.005$, a factor $\approx 45$ below experiment.
$H_5'$ retains the CM anchoring (dynamical preference for the
neighbourhood of $\tau=i$) while accommodating the experimentally
required deviation.
\end{remark}

\begin{axiom*}[$H_7'$, primary: integer-point Damerell ladder]
For the LMFDB CM newform $\fLM$ (CM by $\Qi$, weight $k=5$, level
$N=4$) and the Chowla--Selberg period
$\OmK = \Gamma(1/4)^{2}/(2\sqrt{2\pi})\approx 2.62205755$ for $\Qi$,
the algebraic parts of the critical $L$-values at INTEGER critical
points satisfy
\begin{equation}
\label{eq:ladder}
  \frac{L(\fLM,m)\;\pi^{4-m}}{\OmK^{4}}
  \;=\;
  \begin{cases}
    1/10  & (m=1),\\
    1/12  & (m=2),\\
    1/24  & (m=3),\\
    1/60  & (m=4).
  \end{cases}
\end{equation}
The Cardy density $\rho_{\mathrm{Cardy}}(c)=c/12$ at the unitary
minimal-model central charges $c=1$ (free boson) and $c=1/2$
(Ising) is reproduced \emph{parameter-free} at $m=2$ and $m=3$
respectively.
\end{axiom*}

\begin{axiom*}[$H_7'$, complement: BDP anti-cyclotomic at $s=k/2$]
At the half-integer central point $s=k/2$ outside the
Damerell/Shimura/Deligne archimedean algebraicity domain, the
appropriate algebraic bridge is the anti-cyclotomic $p$-adic
$L$-function of Bertolini--Darmon--Prasanna~\cite{BDP13} for the
CM-by-imaginary-quadratic form.  The $p$-adic value at $s=5/2$ is
algebraic in the Iwasawa algebra of the anti-cyclotomic
$\mathbb{Z}_p$-extension of $\Qi$; we postulate (subject to A1
lit-check) a $p$-adic Cardy bridge for those CFT central charges that
are \emph{not} cleanly fit by the integer ladder of
\eqref{eq:ladder}, viz.\ $c\in\{7/10,4/5\}$.
\end{axiom*}

%%--------------------------------------------------------------------
\section{Main numerical result and the Cardy hits}
\label{sec:main}

We collect the H7-A computation in tabular form
(Table~\ref{tab:ladder}).  All values were obtained by mpmath
($\mathrm{dps}=60$) evaluation of the approximate functional
equation~\cite[Thm.~5.3]{IwaKow04} for $L(\fLM,s)$, with the
sanity-check $L(\fLM,5/2)=0.5200744676\ldots$ matching the LMFDB
record exactly.  The Chowla--Selberg period is
$\OmK = \Gamma(1/4)^{2}/(2\sqrt{2\pi}) = 2.62205755\ldots$, the
classical Lemniscate constant for $\Qi$~\cite{ChowSel49}.

\begin{table}[ht]
\centering
\begin{tabular}{cllcl}
\toprule
$m$ & $L(\fLM,m)$ & $L(\fLM,m)\,\pi^{4-m}/\OmK^{4}$ &
implied $c=12\alpha_m$ & CFT identification \\
\midrule
$1$ & $0.15244713359\ldots$ & $1/10$ & $6/5$ &
exotic, non-unitary minimal\\
$2$ & $0.39910566245\ldots$ & $\bm{1/12}$ & $\bm{1}$ &
\textbf{free boson (Tonks--Girardeau)} \\
$3$ & $0.62691370859\ldots$ & $\bm{1/24}$ & $\bm{1/2}$ &
\textbf{Ising} \\
$4$ & $0.78780300053\ldots$ & $1/60$ & $1/5$ &
no clean unitary minimal CFT \\
\bottomrule
\end{tabular}
\caption{H7-A integer-point Damerell ladder for $\fLM$.  Algebraic
parts $\alpha_m := L(\fLM,m)\,\pi^{4-m}/\OmK^{4}$ are RATIONAL with
denominators $\{10,12,24,60\}$; the boxed entries reproduce
$\rho_{\mathrm{Cardy}}=c/12$ at $c=1$ and $c=1/2$.
The ladder is \emph{standard}: $\alpha_1 = 1/10$ is the
lemniscatic Hurwitz number $H_4$ (Hurwitz 1899~\cite{Hurwitz1899},
prototype example in Harder--Schappacher 1986~\cite{HS86}~\S 1);
$\alpha_2 = 1/12 = B_2/2 = -\zeta(-1)$ is the Bernoulli value
transported to $\Qi$ via Eisenstein--Kronecker / Chowla--Selberg
inside the Damerell--Shimura framework~\cite{Damerell71,Shimura76,ChowSel49}.
The functional equation $L(\fLM,s) \leftrightarrow L(\fLM,5-s)$
forces $\alpha_m/\alpha_{5-m} = \Gamma(5-m)/\Gamma(m)$, hence
$(\alpha_3,\alpha_4) = (\alpha_2/2, \alpha_1/6)$ are
\emph{not independent}; only $(\alpha_1,\alpha_2)$ are.}
\label{tab:ladder}
\end{table}

\begin{theorem}[Bernoulli-anchored Cardy hit]
\label{thm:hit}
Within the Damerell--Shimura algebraicity framework, the algebraic
part $\alpha_2 := L(\fLM,2)\,\pi^{2}/\OmK^{4}$ for $\fLM = $~LMFDB
\texttt{4.5.b.a} (CM by $\Qi$, weight $5$) equals
\begin{equation}
  \alpha_2 \;=\; \tfrac{1}{12} \;=\; \tfrac{B_2}{2} \;=\; -\zeta(-1),
\end{equation}
yielding the Cardy density $\rho_{\mathrm{Cardy}}(1) = c/12 = 1/12$
of the $c=1$ free boson \emph{parameter-free}, where $B_2 = 1/6$ is
the second Bernoulli number.  The functional equation
$L(\fLM,s) = \varepsilon\,L(\fLM,5-s)$ then forces the companion
$\alpha_3 = \Gamma(2)/\Gamma(3)\cdot\alpha_2 = 1/24$, which
reproduces the $c=1/2$ Ising central charge as a structural
$\Gamma$-shadow of the same Bernoulli value.
\end{theorem}

\begin{proof}
The Bernoulli identification $\alpha_2 = B_2/2$ for the lemniscatic
CM-by-$\Qi$ family of weight $\geq 4$ is the classical Eisenstein--%
Kronecker computation: at the integer critical point $m = k-3 = 2$,
the Damerell formula for a CM newform of weight $k$ over an imaginary
quadratic field $K$ collapses to a Bernoulli rational times
$\OmK^{2k-2}/\pi^{k-m-1}$; for $K = \Qi$ and $k=5$ this is
$B_2/2 \cdot \OmK^{4}/\pi^{2}$~\cite{Damerell71,Shimura76,HS86}.
The companion $\alpha_3 = 1/24$ then follows from the
$\Gamma$-factor functional equation, with no independent input.
\end{proof}

\begin{remark}[On the apparent ``two parameter-free hits'']
The empirical observation in Table~\ref{tab:ladder} shows two
Cardy-target rows ($c=1$ at $m=2$ and $c=1/2$ at $m=3$).  These are
\emph{not} two independent algebraic coincidences: the
$\Gamma$-factor functional equation \emph{forces} $\alpha_3 / \alpha_2
= \Gamma(2)/\Gamma(3) = 1/2$, so any $\alpha_2$ yielding a clean
$c$-value at $m=2$ automatically produces a $c/2$ companion at $m=3$.
The genuine parameter-free coincidence is the $c=1$ hit alone --- a
single Bernoulli identity $B_2/2 = 1/12$, with the Ising central
charge appearing as its structural $\Gamma$-shadow.  The strength of
the bridge is correspondingly halved relative to the v7.3 prospectus,
but it remains anchored on a textbook Eisenstein--Kronecker value.
\end{remark}

\begin{remark}[Two-step normalisation]
The structure of \eqref{eq:ladder} factorises as a Damerell-style
binomial $\OmK^{4}/\pi^{4-m}$ times a small rational depending on $m$.
The $\OmK^{4}$ exponent is the $k$-th power expected in
Shimura~\cite{Shimura76} for a weight-$k$ CM form ($k=5$ here, but
the central $L$-values use $\OmK^{2(s-1)}$ with $s$ critical;
$s=m\in\{1,2,3,4\}$ falls in the strip $1\le s\le k-1$).
The numerator pattern $(1/10,1/12,1/24,1/60)$ is what requires
literature anchoring -- see Section~\ref{sec:open}.
\end{remark}

%%--------------------------------------------------------------------
\section{Empirical support for $H_5'$: the W1 attractor scan}
\label{sec:W1}

To test whether $\tau$ near $i$ is dynamically preferred -- as opposed
to merely allowed -- we ran a $30\times 30$ $\chi^{2}$ scan on the
domain $\mathrm{Re}\,\tau\in[-0.5,0.5]$,
$\mathrm{Im}\,\tau\in[0.7,1.5]$ for seven fermion observables
($m_c/m_t$, $m_u/m_c$, $m_e/m_\mu$, $m_\mu/m_\tau$,
$\sin\theta_{12}$, $\sin\theta_{13}$, $\sin\theta_{23}$, all PDG
2024 central values).  The scan ran on a single RTX~5060~Ti GPU
(cosmopower-jax, 386 predictions/s) for $107$~minutes.

\begin{table}[ht]
\centering
\begin{tabular}{lcc}
\toprule
& Best $\tau$ & $|\tau-i|$ \\
\midrule
LYD20 published~\cite{LYD20} & $-0.21+1.52\,i$ & $0.56$ \\
W1 near-$i$ minimum & $-0.1897+1.0034\,i$ & $\bm{0.19}$ \\
\bottomrule
\end{tabular}
\caption{The W1 scan (PC compute, 2026-05-05) finds a local minimum
$\tau^{*}=-0.1897+1.0034\,i$ with $\chi^{2}/\mathrm{dof}=1.05$
(seven observables), a factor $\approx 3$ closer to the
$S$-fixed-point $\tau=i$ than the LYD20 best fit.}
\label{tab:W1}
\end{table}

The seven observables fit the W1 minimum with $\chi^{2}_{\mathrm{up}}
\sim 10^{-19}$, $\chi^{2}_{\mathrm{lep}}\sim 10^{-18}$, all $\chi^{2}$
contribution coming from the CKM sector ($\chi^{2}_{\mathrm{ckm}} =
7.34$).  This is consistent with $H_5'$:
\emph{the empirical attractor lives at $|\tau-i|\sim 0.2$}, and the
scaling is set by the same factor $\sqrt{\epsilon}\sim 0.063$ that
governs the V2 no-go's near-rank-1 obstruction~\cite{V2NoGo} -- with
$|\tau-i|\sim 0.2$ providing precisely the perturbation needed to
break the $S$-collinearity.

\begin{remark}
The W1 scan does not yet include the CKM phase $\delta_{CP}$ nor the
absolute neutrino mass scale; both are degrees of freedom not
constrained by the seven observables used.  A 64-observable refit
including these will be the v7.4.1 update.
\end{remark}

%%--------------------------------------------------------------------
\section{Open questions and pending verifications}
\label{sec:open}

Two independent open questions remain.

\paragraph{Q1: literature anchoring of the rational ladder ---
\textbf{RESOLVED}.}
The pattern $(1/10,1/12,1/24,1/60)$ in Table~\ref{tab:ladder} is the
\emph{standard} CM-by-$\Qi$ Eisenstein--Kronecker ladder for a
weight-$5$ newform: $\alpha_1 = 1/10$ is the lemniscatic Hurwitz
number $H_4$ first computed in Hurwitz~1899~\cite{Hurwitz1899}
(\S~III) as the prototype of the Deligne-conjecture instances later
formalised by Damerell~\cite{Damerell71} and
Shimura~\cite{Shimura76}; $\alpha_2 = 1/12$ is the second Bernoulli
half-value $B_2/2 = -\zeta(-1)$ transported to $\Qi$ via
Chowla--Selberg~\cite{ChowSel49} and recorded as the
prototype example in Harder--Schappacher~1986~\cite{HS86}~\S 1
(reproduced verbatim for $K=\Qi$ in Visser~2021~\cite{Visser21}).
The functional equation reduces $(\alpha_3,\alpha_4) = (\alpha_2/2,
\alpha_1/6)$ to dependent companions, so the ladder has only two
independent rationals, both classical.  $H_7'$(Damerell--ladder) is
therefore a \emph{theorem}, not a conjecture, with the strength
re-stated as in Theorem~\ref{thm:hit} and the Remark following it.

\paragraph{Q2: the unfit central charges $c=7/10$ and $c=4/5$.}
The integer ladder \eqref{eq:ladder} cleanly hits
$c\in\{6/5,1,1/2,1/5\}$ but \emph{not} the unitary minimal central
charges $c=7/10$ (tricritical Ising) and $c=4/5$ (tetracritical
Potts).  Two scenarios are possible:
\begin{itemize}
  \item\textbf{(a)} The complement axiom $H_7'$(BDP) supplies these
    values via $p$-adic anti-cyclotomic $L$-values at $s=k/2$, with
    different small-prime contributions from the
    Iwasawa-algebra interpolation.
  \item\textbf{(b)} They genuinely require a \emph{second} CM
    newform (the candidate would be the level-16 form 16.5.c.a, which
    appears in companion P-NT identifications of hatted
    weight-5 multiplets); in that case the v7.4 anchor must be
    extended to a two-form anchor.
\end{itemize}
Discriminating (a) from (b) is the next concrete test of v7.4.

\paragraph{Q3: Perlmutter degree-4 backup.}
Whether the Rankin--Selberg square $L(\fLM\otimes\fLM,s)$ at
INTEGER critical points $s\in\{5,\ldots,9\}$ supplies the same Cardy
hits as the integer Damerell ladder is the H7-B research question;
both routes should agree on $c=1$ and $c=1/2$ if either is the
correct bridge.

%%--------------------------------------------------------------------
\section{Falsifiers}
\label{sec:falsifiers}

The v7.4 axiomatic restart, like its predecessors, must be
falsifiable.  Three concrete falsifiers are recorded.

\begin{enumerate}
  \item \textbf{Tonks--Girardeau density measurement.}  The Cardy
    formula at $c=1$ predicts a one-loop free-boson density that, in
    the cold-atom Tonks--Girardeau limit, can be measured at the
    percent level.  A deviation from $\rho=1/12$ at the predicted
    operator at the 5\% level falsifies $H_7'$(primary).
  \item \textbf{Modular flavour lepton sector.}  The W1 minimum
    $\tau^{*}=-0.1897+1.0034\,i$ predicts the full 12-parameter
    LYD20 lepton sector (charged-lepton masses, PMNS angles, Majorana
    phases, $m_{\beta\beta}$).  An LYD20-style global fit to PDG
    + KamLAND-Zen + KATRIN + CUPID-Mo-1t at the v7.4.1 stage that
    rejects $|\tau-i|\le 0.3$ at $\ge 5\sigma$ falsifies $H_5'$.
  \item \textbf{NMC cosmological NULL at Cassini.}  An independent
    branch of the ECI programme makes a non-minimally-coupled (NMC)
    scalar-field prediction for the $\gamma$-parameter of the Cassini
    Doppler tracking experiment (PPN bound $|\gamma-1|\le
    2.3\times 10^{-5}$).  The prediction is a NULL: $|\gamma-1|<
    10^{-7}$ at present sensitivity, vanishing in the limit
    $\xi\to 0^{+}$.  A future improvement of the Cassini
    bound to $\le 10^{-6}$ that excludes the NMC NULL falsifies
    the cosmological branch of v7.4.
\end{enumerate}

%%--------------------------------------------------------------------
\section{Conclusion}
\label{sec:conclusion}

The v7.4 amendment converts the strict fixed-point axiom $H_5$ into an
\emph{attractor} statement $H_5'$ supported by a $\chi^{2}$ scan
(W1, 2026-05-05), and replaces the half-integer Damerell--CS bridge
$H_7$ -- refuted by the G1.15 audit on parity grounds -- with a
two-part bridge $H_7'$ comprising (i) an INTEGER-point Damerell
ladder (Table~\ref{tab:ladder}) with parameter-free hits on
$\rho_{\mathrm{Cardy}}=c/12$ at $c=1$ and $c=1/2$, and (ii) a BDP
$p$-adic anti-cyclotomic complement at the original central point.
The CM anchor $\fLM$ (LMFDB CM newform 4.5.b.a, CM by $\Qi$) is
preserved across both rescue routes (H7-A, H7-B, H7-C all converge on
this anchor).  Pending the A1 literature check on the ladder
$(1/10,1/12,1/24,1/60)$, $H_7'$(primary) is offered as a conjecture
backed by a 60-digit numerical fit; the structural converging
recommendation of three independent rescue paths places the burden of
proof, in our view, with a literature anchor rather than with a new
derivation.

%%--------------------------------------------------------------------
\section*{Acknowledgements}

KR thanks the H7-A, H7-B, H7-C, and W1 sub-agent runs (Sonnet harness,
agents\_v647, May~2026) and the Opus G1.15 audit thread for the
refutation analysis.  Computational resources: PC RTX~5060~Ti
(cosmopower-jax 386 predictions/s, JAX named\_shape patch).
LMFDB queries: 4.5.b.a $L$-value cached 2026-05-05.

%%--------------------------------------------------------------------
\section*{Data and code availability}

Repository: \url{https://github.com/AIdevsmartdata/crossed-cosmos}.
Snapshot of v7.3 axioms: Zenodo DOI \Zeno (commit 6.0.53, 2026-05-05
06:00 UTC).  H7-A computation: \texttt{notes/eci\_v7\_aspiration/}
\texttt{H7\_RESCUE/SYNTHESIS\_2026-05-05.md}.  W1 verdict log:
\texttt{notes/eci\_v7\_aspiration/W1\_TAU\_NEAR\_I/results/}
\texttt{W1\_verdict\_2026-05-05.log}.

%%--------------------------------------------------------------------
\begin{thebibliography}{99}

\bibitem{V2NoGo}
K.~Remondi\`ere,
\emph{A no-go theorem for the Cabibbo angle in $S'_4$ modular
flavour models at $\tau=i$},
ECI v7 aspiration note V2 (2026-05-04);
companion paper, this volume.

\bibitem{NPP20}
P.~P.~Novichkov, J.~T.~Penedo, S.~T.~Petcov,
\emph{Double Cover of Modular $S_4$ for Flavour Model Building},
Nucl.~Phys.~B \textbf{963} (2021) 115301;
\arXiv{2006.03058}.
%% [VERIFIED 2026-05-05 via arXiv API.]

\bibitem{LYD20}
X.-G.~Liu, C.-Y.~Yao, G.-J.~Ding,
\emph{Modular invariant quark and lepton models in double covering of
$S_4$ modular group},
Phys.~Rev.~D \textbf{103} (2021) 056013;
\arXiv{2006.10722}.
%% [VERIFIED 2026-05-05 via arXiv API.]

\bibitem{dMVP26}
I.~de~Medeiros~Varzielas, M.~Paiva,
\emph{Quark masses and mixing from Modular $S'_4$ with Canonical
K\"ahler Effects}, \arXiv{2604.01422} (2026).
%% [VERIFIED 2026-05-05 via arXiv API.]

\bibitem{DuWang22}
X.~K.~Du, F.~Wang,
\emph{Flavor Structures of Quarks and Leptons from Flipped SU(5) GUT
with $A_4$ Modular Flavor Symmetry},
JHEP \textbf{01} (2023) 036;
\arXiv{2209.08796}.
%% [VERIFIED 2026-05-05 via arXiv API; NOT "Tanimoto-Yamamoto".]

\bibitem{AbbasKhalil22}
M.~Abbas, S.~Khalil,
\emph{Modular $A_4$ Symmetry With Three-Moduli and Flavor Problem},
\arXiv{2212.10666} (2022).
%% [VERIFIED 2026-05-05 via arXiv API; NOT "Ding-King-Liu-Lu".]

\bibitem{Perlmutter25}
E.~Perlmutter,
\emph{An $L$-function Approach to Two-Dimensional Conformal Field Theory},
\arXiv{2509.21672} (Sept.~2025).
%% [VERIFIED 2026-05-05 via arXiv API; sole author E. Perlmutter.]

\bibitem{Damerell71}
R.~M.~Damerell,
\emph{$L$-functions of elliptic curves with complex multiplication, I},
Acta Arith.\ \textbf{17} (1970/71) 287--301;
\emph{II}, Acta Arith.\ \textbf{19} (1971) 311--317.

\bibitem{Shimura76}
G.~Shimura,
\emph{The special values of the zeta functions associated with cusp forms},
Comm.~Pure Appl.~Math.\ \textbf{29} (1976) 783--804.

\bibitem{Deligne79}
P.~Deligne,
\emph{Valeurs de fonctions $L$ et p\'eriodes d'int\'egrales},
in \emph{Automorphic forms, representations and $L$-functions} (Corvallis, 1977),
Proc.\ Sympos.\ Pure Math.\ XXXIII, Part~2, AMS (1979), 313--346.

\bibitem{BDP13}
M.~Bertolini, H.~Darmon, K.~Prasanna,
\emph{Generalized Heegner cycles and $p$-adic Rankin $L$-series},
Duke Math.\ J.\ \textbf{162} (2013) 1033--1148;
DOI \href{https://doi.org/10.1215/00127094-2142056}{10.1215/00127094-2142056}.
%% [VERIFIED 2026-05-05 via Duke / Project Euclid; no arXiv preprint.]

\bibitem{ChowSel49}
S.~Chowla, A.~Selberg,
\emph{On Epstein's zeta-function (I)},
Proc.~Nat.~Acad.~Sci.~U.S.A.\ \textbf{35} (1949) 371--374.

\bibitem{IwaKow04}
H.~Iwaniec, E.~Kowalski,
\emph{Analytic Number Theory},
AMS Colloquium Publ.\ \textbf{53}, AMS (2004).
% Theorem 5.3 (approximate functional equation) used by H7-A's mpmath
% L-value computation.

\bibitem{Hurwitz1899}
A.~Hurwitz,
\emph{\"Uber die Entwicklungskoeffizienten der lemniskatischen
Funktionen},
Math.~Ann.\ \textbf{51} (1899) 196--226.
% Lemniscatic numbers H_n; H_4 is the prototype CM-by-Q(i) instance
% used here as alpha_1 = 1/10.

\bibitem{HS86}
G.~Harder, N.~Schappacher,
\emph{Special values of Hecke L-functions and abelian integrals},
in: \emph{Workshop Bonn 1984}, Lecture Notes in Math.\ \textbf{1111},
Springer (1986) 17--49 (also LNM \textbf{1399}, Zagier mirror).
% Section 1 prototype example: CM-by-Q(i) lemniscate ladder for
% weight-5 newforms; provides the standard reference for both
% alpha_1 = H_4 = 1/10 and alpha_2 = B_2/2 = 1/12.

\bibitem{Visser21}
R.~Visser,
\emph{Special values of L-functions of CM modular forms},
University of Warwick first-year PhD project (2021).
% Tables for K = Q(sqrt(-7)) and K = Q(i) verifying the Damerell
% ladder structure used in Theorem 2.

\bibitem{ECIv6053}
K.~Remondi\`ere,
\emph{ECI/crossed-cosmos repository, version 6.0.53},
Zenodo (2026-05-05);
DOI \Zeno.
% Snapshot containing v7.3 axioms and the G1.15 refutation log.

\end{thebibliography}

\end{document}
